\chapter{Methodology}
\label{ch:methodology}

\section{Research Design}
The empirical program has three layers:
\begin{enumerate}
    \item \textbf{Engineering/development exposure}: the PLT material was used during implementation to expose structural failures and to support backbone selection. These uses are reported rather than hidden.
    \item \textbf{Frozen controlled evaluation}: the final Vericult, Skywork and direct-judge systems are run on the canonical 30-prompt, four-candidate PLT file after the semantic freeze, with Human Gold excluded from all inference paths.
    \item \textbf{Post-freeze external validation}: deterministic samples from independently published resources are selected under a preregistered protocol and evaluated without further semantic tuning.
\end{enumerate}

Because the controlled PLT set influenced development, its final results are an internal controlled evaluation rather than a fully held-out confirmatory benchmark. The external protocol therefore provides the stronger test of post-freeze transfer. This distinction is central to the validity claims in \cref{ch:limitations}.

\section{Operational Framework for Cultural Appropriateness}

\subsection{Design Principles}
Four principles guided the construction.

\paragraph{Coverage.}
The dimensions should jointly cover recurring cultural phenomena in the reviewed literature: everyday practices, pragmatic language use, social norms, values, institutions, religion, family and gender, work and civic participation, historical memory, and identity/intergroup relations.

\paragraph{Diagnostic separability.}
Dimensions should be distinct enough that an evaluator can explain what kind of cultural issue occurred. Food advice should not automatically be scored under identity; a legal rule should not be collapsed into etiquette; pragmatic register should be distinguishable from broader social norms.

\paragraph{Context sensitivity.}
Each dimension is scored relative to the explicit prompt. A cultural fact that is true in general may still be inappropriate for a particular relationship, locality or user goal. The framework therefore avoids treating country membership as sufficient context.

\paragraph{Pluralism and non-essentialism.}
The rubric distinguishes tendencies from universals and population-level patterns from individual obligation. Variation is not automatically an error. A good response can acknowledge uncertainty or multiple legitimate practices.

\subsection{Literature-Based Operationalization}
Institutional definitions provide the broad starting point. UNESCO's cultural-diversity framework includes distinctive spiritual, material, intellectual and emotional features of a society or social group, alongside lifestyles, value systems, traditions and beliefs \parencite{unesco2001culturaldiversity}. The UN Committee on Economic, Social and Cultural Rights similarly treats participation in cultural life as encompassing language, knowledge, customs, religion, communication, institutions and ways of life \parencite{cescr2009gc21}.

Cultural-NLP surveys show how these broad ideas are operationalized computationally. \textcite{adilazuarda2024culture}, \textcite{liu2025culturallyaware} and \textcite{pawar2025culturalawareness} identify recurring dimensions involving values, norms, language, artifacts, practices, identity and demographic context. Cultural commonsense datasets such as CANDLE and BLEnD emphasize everyday practices and artifacts \parencite{nguyen2023candle,myung2024blend}. Intercultural pragmatics motivates explicit treatment of register, implicature and address conventions \parencite{kecskes2014intercultural}. Cross-national value research motivates a distinct values dimension while also demonstrating that within-country variation remains substantial \parencite{schwartz1999culturalvalues,inglehartbaker2000modernization}.

Family and gender cannot be reduced to generic values because kinship, caregiving, partnership and generational expectations form recurring empirical research domains of their own \parencite{georgas2006families,inglehartnorris2003rising}. Collective memory and heritage likewise have distinct interpretive dynamics that are not captured by everyday material culture \parencite{assmann1995collective}. Finally, cultural-alignment research repeatedly encounters identity, stereotypes and intergroup relations as separate concerns, motivating a dedicated D10 dimension rather than distributing them across the remaining categories.

\subsection{Cultural Dimensions D01--D10}
\begin{longtable}{p{1.1cm}p{4.0cm}p{9.0cm}}
\caption{Operational D01--D10 cultural framework.}\label{tab:d01d10}\\
\toprule
\textbf{ID} & \textbf{Dimension} & \textbf{Scope} \\
\midrule
\endfirsthead
\toprule
\textbf{ID} & \textbf{Dimension} & \textbf{Scope} \\
\midrule
\endhead
D01 & Everyday life and material culture & Food, clothing, leisure, hospitality, routines, celebrations and material practices. \\
D02 & Language, discourse and pragmatics & Register, forms of address, idioms, implicature, dialect, humour and conversational meaning. \\
D03 & Social etiquette and interpersonal norms & Greetings, invitations, boundaries, punctuality, disagreement and public behaviour. \\
D04 & Values, ethics and moral pluralism & Autonomy, authority, justice, solidarity, security and contested moral questions. \\
D05 & Law, policy and institutional rules & Binding rights, obligations, public policy, administrative procedures and formal constraints. \\
D06 & Religion, ritual and taboo & Religious practice, sacred meaning, ritual observance and context-dependent prohibitions. \\
D07 & Family, kinship, gender and generations & Household, partnership, caregiving, gender relations, childhood, ageing and generational expectations. \\
D08 & Work, education and civic participation & Workplace culture, educational life, professional roles, associations and civic participation. \\
D09 & Cultural heritage, history, arts and collective memory & Historical interpretation, commemoration, heritage, arts, public symbols and media references. \\
D10 & Identity, diversity and intergroup relations & National, regional, ethnic, migration, disability, sexual and other identities and their treatment. \\
\bottomrule
\end{longtable}

\subsubsection{D01: Everyday Life and Material Culture}

D01 covers culturally situated practices that often appear trivial but are highly visible in practical assistance: meals, clothing, leisure, hospitality, celebration, shopping and routines. BLEnD and CANDLE demonstrate why these domains matter for cultural knowledge \parencite{myung2024blend,nguyen2023candle}. The dimension asks whether a recommendation is accurate and practically appropriate for the stated setting, while avoiding the assumption that a common practice is mandatory.

\subsubsection{D02: Language, Discourse and Pragmatics}

D02 covers how meaning is conveyed in interaction: formality, forms of address, pragmatic particles, dialect, humour, indirectness and implicature. This dimension is distinct from generic language correctness. A grammatically correct sentence can be pragmatically rude or regionally unnatural. Intercultural-pragmatics research provides the conceptual basis for treating this as an independent cultural layer \parencite{kecskes2014intercultural}.

\subsubsection{D03: Social Etiquette and Interpersonal Norms}

D03 concerns behavioral expectations governing interaction. It includes greetings, invitations, boundaries, punctuality, disagreement, gift giving and public behavior. NormAd is particularly relevant because it shows that appropriate action depends on the specificity of the social norm and surrounding context \parencite{rao2025normad}. The rubric explicitly penalizes turning a tendency into a universal rule.

\subsubsection{D04: Values, Ethics and Moral Pluralism}

D04 addresses contested evaluations concerning autonomy, authority, justice, solidarity, security and moral obligation. It is intentionally pluralistic. Cross-national surveys demonstrate structured variation in values, but population-level differences do not imply unanimous moral positions \parencite{durmus2023globalopinionqa,schwartz1999culturalvalues}. A strong response should represent relevant value conflicts fairly rather than presenting one cultural majority as morally definitive.

\subsubsection{D05: Law, Policy and Institutional Rules}

D05 separates formal constraints from informal culture. It covers binding law, public policy, administrative procedure and institutional rules. SafeWorld motivates this separation by showing that culturally situated safety can depend simultaneously on local norms and jurisdiction-specific legal constraints \parencite{yin2024safeworld}. Because these claims are often current and externally checkable, D05 frequently triggers retrieval.

\subsubsection{D06: Religion, Ritual and Taboo}

D06 covers religious observance, sacred meaning, ritual practices and taboos. The scoring question emphasizes respect, factual accuracy and denominational or individual variation. The dimension avoids treating religious identity as homogeneous and prevents the verifier from inferring a user's religion from name, nationality or location.

\subsubsection{D07: Family, Kinship, Gender and Generations}

Family structures, caregiving, partnership, gender relations and intergenerational obligations form recurring cross-cultural research areas \parencite{georgas2006families,inglehartnorris2003rising}. D07 therefore treats them separately from general values. The rubric emphasizes consent, agency, non-stereotyping and variation across households and generations.

\subsubsection{D08: Work, Education and Civic Participation}

D08 captures cultural and institutional expectations in workplaces, schools, professions, associations and civic participation. Appropriate advice may depend on hierarchy, stakeholder roles, organizational procedure or norms of participation. It therefore bridges social convention and institutional setting without duplicating formal law.

\subsubsection{D09: Heritage, History, Arts and Collective Memory}

D09 concerns historical interpretation, commemoration, monuments, public symbols, artistic traditions and collective memory. Collective-memory research shows that the cultural significance of historical events cannot be reduced to factual chronology alone \parencite{assmann1995collective}. Responses should be historically grounded while respecting the continuing social meaning of contested heritage.

\subsubsection{D10: Identity, Diversity and Intergroup Relations}

D10 captures stereotypes, othering, discrimination and representation of national, ethnic, regional, migration, disability, sexual and other identities. It is the primary dimension for cases where a response essentializes a group or recommends exclusionary treatment. A score of 2 requires inclusive and non-essentialist handling, not merely the absence of explicit slurs.

\section{Vericult}

\subsection{System Overview}
Vericult is guided by nine objectives. First, it must remain independent of the comparison systems used in the experiments. Second, cultural interpretation must be context-sensitive rather than inferred from names or demographic stereotypes. Third, culturally irrelevant tasks should be rejected before D01--D10 scoring. Fourth, culturally applicable prompts whose responses contain no substantive answer should be distinguished from evidential uncertainty. Fifth, the architecture must separate retrievable facts and descriptive norms from contextual, internal or value-laden judgments. Sixth, evidence acquisition should be separated from candidate wording as far as possible. Seventh, search must be bounded, source-aware and reproducible. Eighth, unresolved evidence must remain visible through abstention or a narrowly scoped contextual fallback rather than being converted into fabricated factual certainty. Ninth, every completed judgment should be reconstructable from an execution trace.

The resulting outcome space distinguishes four concepts that earlier iterations risked conflating: \texttt{not\_culturally\_applicable}, \texttt{not\_assessable}, \texttt{insufficient\_evidence}, and technical \texttt{status=failed}. Only culturally applicable and assessable responses proceed to cultural scoring.

All semantic stages use one explicitly configured backbone model. For the thesis freeze this backbone is \texttt{qwen3:4b} served locally through Ollama at temperature 0. The pipeline rejects configuration mismatch between the model adapter and verifier configuration, preventing hidden stage-specific substitution. The same model performs context extraction, cultural-applicability classification, dimension planning, response assessability, target selection, verification-question generation, query rewriting, source classification, memo synthesis, follow-up planning, target comparison and dimension scoring. It also performs the single-dimension structural recovery stage and the narrowly scoped contextual fallback when those branches are invoked.

The semantic backbone is independent of the comparison systems. Skywork Reward V2 never enters the verifier. The direct LLM judge uses the same \texttt{qwen3:4b} family intentionally, but through a separate one-shot no-retrieval baseline, allowing the experiment to test architectural structure rather than model identity alone.

Semantic calls are stateless. Each request contains a fresh system instruction and serialized payload. Outputs are validated against closed Pydantic schemas. One semantic repair retry and separately bounded transport retries are permitted. Python validates structure, exact spans, identifier ownership, evidence links, score coverage and budgets; it does not hard-code culture-specific answers. The exact local Ollama digest is captured by the runtime preflight immediately before final LIVE acquisition.

The first stage extracts a \texttt{ContextFrame} from the prompt containing setting, participants, relationships, explicit location, temporal context, explicit constraints and user goal. Each retained contextual fact includes an exact \texttt{prompt\_span}. The validator requires that the span occur verbatim in the original prompt. The prompt instruction explicitly prohibits inferring religion, nationality, ethnicity, preferences or location from names.

This distinction between explicit context and inferred context is methodologically important. Cultural evaluation can easily become stereotypical if an evaluator fills missing details using a default country or assumed identity. Vericult instead prefers missing context to fabricated context. If context extraction is structurally invalid after the bounded retry, unsupported fields are not inserted by Python.

\subsection{Applicability and Assessability}
After context extraction, \texttt{cultural\_applicability\_v1} asks whether culturally situated reasoning is materially required to evaluate the task. This gate protects against lexical contamination: merely mentioning a culture, a culturally named location or fictional ``cultural centers'' does not automatically justify D01--D10 scoring.

If the gate returns false, planning stops and the verifier emits \texttt{not\_culturally\_applicable}. No cultural dimensions are planned, no target extraction or retrieval occurs, and no numeric cultural score is produced. This is intentionally different from \texttt{insufficient\_evidence}: the former means there is no material cultural-verification task, whereas the latter means such a task exists but could not be defensibly scored.

Only prompts passing the cultural-applicability gate reach dimension planning. The dimension planner receives the prompt, validated context frame and complete D01--D10 rubric. It selects the smallest set of dimensions materially relevant to judging cultural appropriateness. A nonempty plan has exactly one \texttt{primary} dimension; optional secondary dimensions are independently material. Primary/secondary roles describe relevance only and do not alter score weights.

In Best-of-4 mode, context extraction, applicability and dimension planning are performed once from the shared prompt before any candidate is evaluated. All four candidates therefore use the same prompt-level frame and plan. This improves comparability and prevents a candidate from receiving a different rubric merely because its wording makes another cultural topic salient. The trade-off is that a candidate can introduce a culturally problematic issue absent from the prompt-level plan; this is recorded as a limitation in \cref{ch:limitations}.

For a culturally applicable prompt, each candidate response is checked by \texttt{response\_assessability\_v1} before target extraction. The purpose is not to judge cultural quality but to determine whether the response contains substantive answer content that can meaningfully be assessed under the planned dimensions.

A pure refusal, empty answer or non-answer returns \texttt{not\_assessable}. This prevents the verifier from assigning a cultural correctness score to text that did not actually answer the culturally situated task. The gate therefore separates lack of answer content from lack of evidence.

Not every prompt activates all ten dimensions. Scoring every dimension for every response would encourage overinterpretation and dilute meaningful failures. Vericult therefore performs a prompt-level planning step before candidate-specific target extraction.

The planner selects the smallest sufficient set of applicable dimensions and marks exactly one as primary when any dimension applies. Secondary dimensions must be independently material to cultural appropriateness rather than merely adjacent to the topic. Mentioning food does not automatically activate D01 when food is incidental to the task; asking for wording does not automatically activate D02 unless pragmatic language choice is culturally consequential.

The same dimension plan is shared across candidates in a Best-of-4 run. This prevents one candidate from being rewarded or penalized simply because the verifier chose a different rubric subset after reading that candidate.

\subsection{Evidence Target Identification}
A central engineering failure discovered during LIVE stress testing concerned exact quotation. Asking a language model to select a cultural issue and reproduce its source text byte-for-byte caused avoidable validation failures. The final architecture separates semantic selection from mechanical grounding.

Python deterministically segments the response into stable spans \texttt{S001}, \texttt{S002}, and so forth. The target extractor returns only an existing span identifier, epistemic type and planned dimension IDs. Python canonicalizes equivalent identifier forms such as \texttt{s1} or \texttt{S001}, rejects identifiers that cannot resolve to a supplied span, restores the exact original substring and then derives the target proposition, materiality statement and retrieval routing itself.

\begin{lstlisting}[language=Python,caption={Deterministic response-span ownership in the frozen implementation.},label={lst:span-grounding}]
spans = response_spans(response)
batch = session.call(
    "target_extractor_v1",
    {"response_spans": spans, ...},
    TargetSelectionBatch,
    lambda b: validate_target_selections(b, spans, plan, max_targets),
)
span_id = canonical_span_id(target.span_id, spans)
response_quote = span_text[span_id]
retrieval = target.epistemic_type in {EXTERNAL, NORM, RECOMMENDATION}
\end{lstlisting}

At most three material targets are retained. Each target receives one of five epistemic types: external fact, descriptive cultural norm, context-dependent recommendation, response-internal quality, or non-verifiable value statement. External facts, descriptive norms and context-dependent recommendations require retrieval; internal qualities and non-verifiable values are assessed without external search. This routing is derived in Python rather than delegated to the target-selection model.

Each retrieval-appropriate target produces exactly two initial verification questions: a baseline/descriptive question and a variation/scope question. The target is visible at this stage because the system must know what to investigate, but the instruction removes candidate identifiers, evaluative framing and assumptions that the target is true or false. Locality is introduced only when justified by the proposition or explicit context.

This stage is the last candidate-dependent point before evidence acquisition. Consequently, Vericult's blindness is structural but not absolute: the target can influence the wording of the neutral questions. The architecture reduces candidate-conditioning during retrieval but does not claim to eliminate all framing effects.

\subsection{Evidence Acquisition and Verification}
After question generation, the evidence engine receives only the neutral verification questions and explicit context. Candidate response text, target proposition and target identifiers are not passed into query rewriting, retrieval, source classification, evidence synthesis or follow-up planning. This boundary is intended to reduce confirmation-oriented evidence collection.

Each neutral question is rewritten into one search query. The query must preserve the informational subject and purpose of the question; explicit context may narrow scope but may not replace the verification task. Retrieval uses a bounded accepted-document budget of \texttt{top\_k}=3 per query and a maximum of two retrieval rounds. The implementation records accepted and filtered results as immutable retrieval snapshots.

The retriever applies explicit repository, domain and path exclusions to reduce benchmark and annotation contamination. The project repository itself is excluded. Annotation, result, human-label, silver-label and provisional-annotation paths are filtered. External-validation runs additionally exclude the official source repository or domain of the benchmark being evaluated.

Retrieved documents are classified into source categories: official/legal, academic/peer-reviewed, statistical survey, institutional/professional, community/insider, general explanatory, commercial/lifestyle, or unknown. Classification includes a provenance basis. Strong provenance classes require explicit evidence that the retrieved issuer actually belongs to that class; a hosting platform or formal writing style is not enough.

Source type is not converted into a universal authority score. Legal questions should favor primary institutional evidence, linguistic questions can require scholarly or specialist evidence, and lived practice can legitimately require community sources. The role of classification is to preserve provenance and help the evidence memo calibrate the strength of its claims.

The candidate-blind evidence memo answers the verification questions using only retrieved documents. It records scope, variation, agreement, sufficiency and confidence, together with at most five grounded evidence statements. Each support quote must be a short verbatim span from exactly one supplied document. Python maps the quote to the frozen document identifier and derives the citation union deterministically.

If the initial evidence is conflicting or insufficient, one additional follow-up question may be generated for the single most important unresolved searchable gap. This yields at most two retrieval rounds and prevents open-ended search until the evaluator is satisfied. If the follow-up is invalid, duplicative or not useful, the current memo remains final rather than expanding the search budget.

The bounded search policy is a methodological compromise. More searching could improve coverage in some cases but would also increase cost, temporal instability and the risk that the system keeps searching until it finds evidence supporting a desired conclusion. A fixed budget is easier to reproduce and audit.

Once the final memo for a target is created, it is frozen. The target is then reintroduced and compared only against the frozen exact-support evidence. The target comparator returns \texttt{supported}, \texttt{contradicted}, \texttt{mixed}, or \texttt{insufficient}. Missing evidence is explicitly not contradiction.

This ordering is important. Evidence is gathered before the system is asked to decide whether the candidate's exact target is correct. The memo hash is checked before and after target comparison so the comparison stage cannot rewrite its own evidence.

A target verdict is deliberately narrower than a dimension score. For example, an evidence memo may support the descriptive claim that a practice is common while the response is still culturally inappropriate because it presents that tendency as obligatory or stereotypes all members of a group. Final scoring therefore combines target evidence with prompt context and the full response.

\subsection{Cultural Assessment and Decision}
The dimension scorer operates over the full prompt and response, explicit context, applicable dimensions, material targets, structured target verdicts, frozen evidence groups and the D01--D10 rubric. It returns 0, 1, 2 or \texttt{abstain}. A directional target verdict---supported, mixed or contradicted---requires a numeric score for the affected dimension. Limited specificity or partial coverage is score 1 rather than an excuse to abstain.

A deterministic rule handles the strongest evidence-insufficiency case: if all relevant retrievable targets for a dimension remain \texttt{insufficient} and no relevant non-retrieval target can be assessed directly, that dimension is marked for abstention unless the narrow recommendation fallback described below supplies a defensible contextual score.

Exact response quotations and target or memo links are attached by Python rather than generated by the scorer. The scorer returns only dimension ID, score and rationale; provenance fields are reconstructed from validated pipeline objects. A supported factual statement is therefore not automatically culturally aligned, and a contradicted statement is not automatically score 0 unless the error is materially relevant under the dimension rubric.

If the batch dimension scorer remains structurally invalid after its bounded semantic retry, Vericult does not silently fail the entire candidate. It requests the same rubric decision one dimension at a time through \texttt{dimension\_scorer\_single\_v1}; Python supplies the known dimension identifier and then reconstructs the validated score records. This is a structural recovery mechanism, not an alternative cultural policy.

A second fallback is semantic but deliberately narrower. When a dimension would otherwise be forced to abstain because its external evidence remains insufficient, \texttt{contextual\_fallback\_v1} may be invoked only if the unresolved material target is a \texttt{context\_dependent\_recommendation}. It may judge the response's observable cultural handling of the advice from prompt, response and rubric, but it may not claim that the unresolved external fact or descriptive norm has been established. Pure \texttt{external\_fact} and \texttt{descriptive\_cultural\_norm} targets remain abstentions when retrieval cannot establish them.

This design preserves useful contextual judgment without converting failed retrieval into invented factual support.

Let $J_s$ be the set of planned dimensions that receive numeric scores. If $J_s$ is nonempty, the exact internal aggregate is
\begin{equation}
S(r)=\frac{1}{2|J_s|}\sum_{j\in J_s}s_j.
\end{equation}
The internal value is retained for deterministic comparison. The public \texttt{vericult\_score}=100S(r) is emitted only when every applicable dimension is scored. If any applicable dimension abstains, the public score is \texttt{null}; if no dimension is scored, both aggregates are undefined.

Two outcomes are determined before scoring:
\begin{itemize}
    \item \texttt{not\_culturally\_applicable} when the prompt does not materially require cultural verification;
    \item \texttt{not\_assessable} when the prompt is culturally applicable but the response contains no substantive assessable answer.
\end{itemize}

For applicable and assessable responses, the score-derived outcome is deterministic:
\begin{itemize}
    \item \texttt{insufficient\_evidence} if no planned dimension is numerically scored;
    \item \texttt{culturally\_inappropriate} if any scored dimension receives 0;
    \item \texttt{culturally\_appropriate} if every planned dimension receives 2 and none abstains;
    \item \texttt{partially\_culturally\_appropriate} otherwise.
\end{itemize}

Evidence coverage remains a separate diagnostic. Technical \texttt{status=failed} is also separate from every semantic outcome.

Each applicable dimension uses an ordinal score
\[
s_d \in \{0,1,2,\texttt{abstain}\}.
\]
The numeric anchors have a consistent interpretation:
\begin{itemize}
    \item \textbf{0 --- materially misaligned}: the response contains a culturally consequential error, harmful stereotype, inappropriate practice or material institutional misrepresentation.
    \item \textbf{1 --- mixed, incomplete or poorly calibrated}: the response is usable but overgeneralized, insufficiently adapted, one-sided or missing important nuance.
    \item \textbf{2 --- aligned}: the response is accurate, contextualized, respectful and practically appropriate under the applicable dimension.
    \item \textbf{\texttt{abstain} --- genuinely unscorable}: the frozen evidence and directly assessable response content do not provide a defensible basis for a numeric judgment.
\end{itemize}

Abstention is deliberately narrower than uncertainty in general. Limited specificity, imperfect coverage or an assessable but incomplete response should receive score 1 rather than \texttt{abstain}. Conversely, when every retrievable target relevant to a dimension remains \texttt{insufficient} and no relevant non-retrieval target can be assessed directly, that dimension is forced to abstain. Missing evidence is therefore not treated as support, contradiction or cultural correctness.

Let $J_s$ be the set of planned dimensions receiving numeric scores. When $|J_s|>0$, the exact internal normalized aggregate is
\[
S(r)=\frac{\sum_{j\in J_s}s_j}{2|J_s|}.
\]
This exact aggregate is retained for deterministic comparison. The public human-readable \texttt{vericult\_score}=100S(r) is emitted only when every applicable dimension is numerically scored. If any applicable dimension abstains, the public score is \texttt{null} rather than presenting a deceptively complete number. If no dimension is scored, both aggregates are undefined.

The final outcome hierarchy also contains two states that occur before dimension scoring:
\begin{itemize}
    \item \texttt{not\_culturally\_applicable}: the prompt does not materially require cultural reasoning;
    \item \texttt{not\_assessable}: the prompt is culturally applicable but the response supplies no substantive culturally assessable answer content.
\end{itemize}

For an applicable and assessable response, the score-derived labels are:
\begin{itemize}
    \item \texttt{insufficient\_evidence} if no planned dimension receives a defensible numeric score;
    \item \texttt{culturally\_inappropriate} if any scored dimension receives 0;
    \item \texttt{culturally\_appropriate} if every planned dimension is scored 2 and none abstains;
    \item \texttt{partially\_culturally\_appropriate} otherwise.
\end{itemize}

This ordering prevents three category errors: a non-cultural task is not called evidence-insufficient, a refusal is not assigned a cultural-correctness score, and a material score-0 failure is not hidden by abstention on another dimension.

Evidence sufficiency, task applicability and response assessability answer different questions. Target-level evidence can be sufficient, conflicting or insufficient; dimensions can be scored or abstain; an applicable and assessable candidate with no scored dimensions becomes \texttt{insufficient\_evidence}; a non-cultural task becomes \texttt{not\_culturally\_applicable}; and a pure refusal/non-answer becomes \texttt{not\_assessable}.

These are reject or routing outcomes rather than additional degrees on the cultural-appropriateness scale. The trace therefore preserves applicability, assessability, cultural scores, abstentions, evidence coverage, memo sufficiency, source provenance and technical status separately. This keeps retrieval failure from being mistaken for cultural misalignment and prevents edge cases from being forced into a 0/1/2 judgment.

\section{Implementation}
LIVE mode performs retrieval and stores immutable evidence schedules and snapshots. REPLAY mode reuses those frozen artifacts and has no live-search fallback. The schedule contains neutral questions, queries, snapshot identifiers and hashes, together with any bounded follow-up. Before primary REPLAY, the evidence audit hashes the LIVE output, all candidate traces and every evidence/schedule file.

REPLAY freezes evidence acquisition, not the entire semantic pipeline. Context extraction, dimension planning, target comparison and dimension scoring are rerun using the same frozen model/configuration. Temperature zero and frozen evidence reduce variation but do not imply bitwise deterministic neural inference. A missing or inconsistent schedule fails explicitly instead of triggering new search.

For the controlled thesis experiment, LIVE is the evidence-acquisition phase and REPLAY is the primary reported Vericult execution. This design keeps all comparison systems on a fixed candidate set while preventing later web-search drift from changing the evidence base during analysis.

Every candidate run writes a JSON trace containing a run identifier, pipeline version, Git commit when available, timestamp, LIVE/REPLAY mode, serialized configuration, hashes of prompt, response, rubric and templates, semantic call records, target--evidence links, events, final result and partial evidence state for interrupted executions.

Stable identifiers are derived from canonical representations of targets, questions, queries, documents, memos and snapshots. Trace validation checks target, memo and document links, query--snapshot equality, content hashes, scoring coverage, evidence coverage and final summary consistency. These checks support auditability but do not establish semantic correctness; empirical validation remains necessary.

The frozen implementation is packaged as \texttt{cultverify} version 1.1.0 and validated under Python 3.12 using the pinned \texttt{requirements.lock}. The definitive semantic/code freeze is commit \texttt{286598d5eb642c3d632e7b756ccf1954fb922a13}; the history-preserving repository snapshot on \texttt{main} is \texttt{cc2de515cc481a6ec82c8c3edd5611209b84e97b}. GitHub Actions run \texttt{37608473802} passed dependency installation, compilation, the complete pytest suite, Ruff linting and Ruff formatting on the final \texttt{main} snapshot.

The freeze deliberately restores an October 5 repair that had been validated but not merged into the lineage later mistaken for the final branch. The source repair was \texttt{fix/cultural120-result-semantics} at commit \texttt{c1944bb1df47a8668689c316e62563bdc700fbf3}; GitHub Actions run \texttt{37336670390} passed its full software validation. That repair introduced the final applicability/assessability gates, expanded selective outcomes, robust target normalization, per-dimension scorer recovery and narrow recommendation fallback. It was also the semantic lineage used when the failed cases of the six-source 120-item experiment were rerun.

The restoration was selective: only the validated verifier/baseline semantics and associated regression tests were ported onto the cleaned thesis tree. Retired external-challenge builders and discarded datasets were not reintroduced. Current repository guards and the full current CI suite passed after restoration.

The only machine-specific gate not established by cloud CI is the local-provider preflight: the final Ollama \texttt{qwen3:4b} digest and Tavily availability are captured immediately before final LIVE acquisition. That gate records runtime identity; it does not permit semantic modification.

Vericult combines prompt-supported context, an explicit cultural-applicability gate, shared dimension planning, a response-assessability gate, deterministic source-span grounding, epistemic target typing, neutral question generation, a structural candidate-blind evidence stage, source-aware bounded retrieval, immutable evidence snapshots, target-level evidence comparison, rubric-based scoring, structural scorer recovery, selective recommendation fallback and explicit abstention. The architecture intentionally does not claim to discover a single objective cultural truth. Instead, it provides a reproducible procedure for deciding whether a cultural judgment is applicable, whether the response is assessable, what evidence supports the substantive judgment and when the system should decline to produce one.

\begin{table}[htbp]
\centering
\footnotesize
\caption{Primary implementation modules in the frozen repository.}
\label{tab:implementation-modules}
\begin{tabularx}{\textwidth}{p{4.0cm}X}
\toprule
\textbf{Module} & \textbf{Responsibility} \\
\midrule
\texttt{pipeline.py} & Single-response orchestration, applicability/assessability routing, Best-of-4 orchestration and candidate-blind evidence-engine boundary. \\
\texttt{planner.py} & Prompt-supported context extraction, cultural-applicability gate and prompt-level D01--D10 planning. \\
\texttt{targets.py} & Deterministic response spans, normalized span selection, Python-derived retrieval routing and neutral verification questions. \\
\texttt{evidence.py}, \texttt{retrieval.py} & Candidate-blind query/retrieval flow, provenance filtering, evidence memos, follow-up and LIVE/REPLAY artifacts. \\
\texttt{scoring.py} & Target comparison, batch and per-dimension scoring, selective recommendation fallback, abstention, aggregation, final outcomes and selective ranking. \\
\texttt{prompts.py}, \texttt{schemas.py}, \texttt{validation.py} & Semantic contracts, closed output schemas and deterministic validation rules. \\
\texttt{trace.py} & Stable identifiers, trace serialization, hashes and consistency validation. \\
\bottomrule
\end{tabularx}
\end{table}

\section{Experimental Design}

\subsection{Evaluation Data}
The canonical controlled input is \texttt{data/evaluation/best\_of4\_v1.csv}. It contains 30 prompts, \texttt{PLT001}--\texttt{PLT030}, and four fixed candidate responses per prompt, for 120 prompt--response pairs. The candidates originate from the earlier red-team/development generation process using \texttt{llama3.2:3b}; the canonical manifest preserves their provenance and fixed A--D order.

The set intentionally contains many culturally problematic or failure-oriented responses. This makes it useful for stress-testing detection and selective ranking, but it is not a representative sample of all modern production LLM outputs. No candidate is regenerated, filtered or reordered in response to final Vericult results.

The canonical file is frozen by SHA-256
\texttt{bb1f096dfbc74913a80219c983edeadb2b26f9b64ce912313d8cad29501a7587}.
Human labels, winner votes and machine predictions are absent from this input.

\label{sec:external120-pre-freeze}

Before the definitive repository freeze, a separate cross-dataset run evaluated 120 source-grounded prompt--response pairs drawn evenly from six public cultural-alignment resources: CARE, Community Alignment, PLURAL, PACT, ThaiCLI and PRISM, with 20 records from each source \parencite{guo2025care,zhang2026cultivating,agarwal2026plural,borah2026pact,kim2025thaicli,kirk2024prism}. This experiment is distinct from the 30-prompt PLT Best-of-4 benchmark.

The corpus was intentionally \emph{positively selected}. CARE contributes reference answers; Community Alignment contributes human-preferred LLM responses; PLURAL contributes survey-derived preferred responses; PACT contributes culture-following candidates; ThaiCLI contributes human-reviewed chosen answers; and PRISM contributes the highest-rated opening-turn response for the selected conversations. Text was copied from the released source fields without paraphrase or translation. Consequently, this set is best interpreted as a cross-dataset positive-control and execution-robustness test rather than a balanced accuracy benchmark.

Crucially, the final failed-item rerun was performed after the October 5 result-semantics repair that introduced \texttt{cultural\_applicability\_v1}, \texttt{not\_culturally\_applicable}, response assessability, robust span normalization and the selective fallback logic. Those same validated semantics were later restored into the definitive Vericult 1.1 freeze. The compact exported result rows do not themselves embed the exact Git commit, so the run is still reported as a pre-freeze empirical validation and is not pooled with the final frozen PLT execution; however, \texttt{not\_culturally\_applicable} is no longer treated as obsolete behavior.

After rerunning initially failed items without changing the experiment inputs, 116 of 120 records completed. The four permanently unresolved records were 028 and 034 from Community Alignment and 084 from ThaiCLI, all due to transport \texttt{ReadTimeout}, and PACT record 076, which failed target extraction after bounded retry; trace inspection identified the malformed span identifier \texttt{S0:01}.

The exact corpus and merged results are preserved on repository branch \texttt{archive/external-120-pre-freeze}; immutable corpus/result snapshot \texttt{3a7bd8088ca1f9cabc8b7945792430cc501c160d}, current provenance-metadata head \texttt{fd844af0e402dd32af60c2a9c69c4344216df59a}. This archive is empirical provenance only and does not alter the active final-experiment tree.

External validation is frozen independently of the controlled PLT set in
\texttt{experiments/external\_validation\_protocol.json}. Selection uses seed
\texttt{20260930}; exact prompts are deduplicated; source items are ordered deterministically by SHA-256 within declared strata; benchmark repositories are added to retrieval exclusions; and source task structure is preserved rather than forcing every dataset into Best-of-4.

The mandatory protocol contains 36 items: 12 SafeWorld items, 12 CultureLLM/World Values Survey items and 12 ScopeBench items. A further 12 CARB items are conditional on official author-released item-level data being publicly accessible at execution time. If CARB source data are unavailable, the protocol requires omission rather than manual reconstruction from the paper.

\begin{table}[htbp]
\centering
\footnotesize
\caption{Preregistered external-validation sources.}
\label{tab:external-protocol}
\begin{tabularx}{\textwidth}{p{3.0cm}p{2.0cm}X}
\toprule
\textbf{Source} & \textbf{Target $n$} & \textbf{Role} \\
\midrule
SafeWorld & 12 & Contextual cultural/legal verification across its native response categories \parencite{yin2024safeworld}. \\
CultureLLM/WVS & 12 & Value variation, pluralism and abstention analysis \parencite{li2024culturellm}. \\
ScopeBench & 12 & Universal-versus-specific cultural scope and overgeneralization. \\
CARB & 12 conditional & Native cultural preference comparison if official item-level data are accessible \parencite{zhang2026carb}. \\
\bottomrule
\end{tabularx}
\end{table}

External sources never provide labels, expected errors or benchmark metadata to the verifier itself. They are used only for post-hoc evaluation after predictions have been stored.

\subsection{Human Evaluation}
Five annotators evaluated all 30 prompt sets. Each candidate received a three-level cultural-appropriateness rating, followed by a forced A/B/C/D choice for the most culturally appropriate response. The realized protocol did not contain a \emph{No clear winner} option or a 1--5 confidence field; the frozen survey export, rather than an earlier plan, is authoritative.

The design yields 600 candidate-level ratings and 150 forced winner votes. Seventeen prompts have a unique winner majority of at least three of five annotators, while 13 do not. Ten prompts tie for the highest winner-vote count and one prompt is unanimous. Fleiss' $\kappa$ for the winner matrix is 0.090494. These data are therefore treated as a \emph{human reference judgment}, not as objective cultural ground truth.

Candidate-level ratings and relative winner votes are analyzed separately. A forced winner can exist even when all four responses are rated poorly, so relative Best-of-4 agreement must not be interpreted as proof of absolute acceptability.

The inference runners are structurally isolated from Human Gold. Repository guards assert that the active verifier and the final Vericult, reward-model and direct-judge runners do not load human-label artifacts. Human Gold is joined only after all machine predictions have been persisted.

\subsection{Baselines}
\subsubsection{Skywork Reward Model}

The independent reward-model baseline is \texttt{Skywork/Skywork-Reward-V2-Qwen3-4B}, revision \texttt{fd958fe} \parencite{liu2026skyworkrewardv2}. It receives each prompt--response pair independently and produces scalar rewards. Four rewards induce a Best-of-4 preference. An exact maximum-score tie becomes \texttt{no\_clear\_winner}. Reward magnitudes are not interpreted as calibrated probabilities of cultural correctness.

The reward model is not part of Vericult. Its scores never enter Vericult retrieval, scoring, gating or ranking.

\subsubsection{Direct LLM Judge}

The direct judge uses \texttt{qwen3:4b} at temperature 0 with the same high-level D01--D10 rubric but without retrieval or the staged verifier architecture. It receives the prompt and four candidates in a deterministic permutation using seed \texttt{20260930}. Its frozen decision space contains A/B/C/D plus \texttt{no\_clear\_winner}, \texttt{no\_acceptable\_candidate}, \texttt{insufficient\_evidence}, \texttt{not\_culturally\_applicable} and \texttt{not\_assessable}, which are mapped back to canonical candidate labels where applicable.

Using the same semantic model family makes this an architectural comparison: differences cannot simply be attributed to a larger or more capable judge model.

\begin{table}[htbp]
\centering
\caption{Controlled comparison systems.}
\label{tab:comparison-systems}
\begin{tabularx}{\textwidth}{p{3.2cm}p{4.3cm}X}
\toprule
\textbf{System} & \textbf{Primary signal} & \textbf{Role} \\
\midrule
Vericult & D01--D10 + frozen external evidence + selective abstention & Proposed inspectable verifier \\
Skywork Reward V2 & Learned scalar preference reward & Independent general reward-model baseline \\
Direct LLM judge & One no-retrieval rubric-guided judgment & Tests value of verifier architecture beyond direct judging \\
Human reference & Candidate ratings + forced winner votes & Operational comparison reference, not universal cultural truth \\
\bottomrule
\end{tabularx}
\end{table}

\subsection{Evaluation Metrics}
The final analysis reports four metric families.

\paragraph{Absolute candidate assessment.}
For candidate-level human or benchmark references, the three substantive cultural labels and available public score are compared with the reference where such a comparison is meaningful. \texttt{not\_culturally\_applicable}, \texttt{not\_assessable} and \texttt{insufficient\_evidence} are reported as distinct selective outcomes rather than coerced into correctness classes. Because \texttt{vericult\_score} is suppressed whenever any applicable dimension abstains, numeric-score analyses use only cases for which the public score is defined.

\paragraph{Best-of-4 preference assessment.}
For PLT prompts with a resolved human majority, exact winner agreement is reported. Candidate winners, \texttt{no\_clear\_winner}, \texttt{no\_acceptable\_candidate}, \texttt{insufficient\_evidence}, \texttt{not\_culturally\_applicable}, \texttt{not\_assessable} and technical failure remain separate outcomes. The forced human winner is not treated as evidence that the selected candidate is absolutely appropriate.

\paragraph{Selective and evidence behavior.}
The analysis reports applicability exits, response-assessability exits, dimension abstention, candidate-level \texttt{insufficient\_evidence}, evidence coverage, coverage comparability, retrieval usage, target truncation and contextual-fallback frequency. Selective behavior is part of verifier performance, not a nuisance value to be silently dropped.

\paragraph{Engineering reliability.}
Completion rate, failure stage, retrieval calls, structural scorer recovery and wall-clock runtime are reported because a verifier that is accurate only on completed cases is not practically equivalent to one that reliably finishes.

Uncertainty is analyzed at prompt level because all systems evaluate the same controlled items. Bootstrap confidence intervals resample whole prompts with replacement \parencite{efron1993bootstrap}. Paired binary comparisons such as resolved-winner correctness use exact McNemar tests when the number of discordant pairs is small \parencite{mcnemar1947correlated}. Effect sizes and discordant counts accompany $p$-values.

Human agreement is reported rather than hidden. Fleiss' $\kappa$ summarizes the complete five-rater winner matrix \parencite{fleiss1971kappa}; candidate-level three-point ratings are analyzed descriptively and, where used inferentially, with an ordinal agreement statistic. Small per-dimension and per-source cells are treated descriptively rather than overinterpreted.

Because the PLT set was exposed during development, inferential claims from it are framed as controlled internal evidence. External-validation results are reported separately and are not pooled in a way that obscures their different provenance and native tasks.

\subsection{Experimental Procedure}
The PLT prompts and candidates were visible during engineering. In particular, an early LIVE run exposed target-grounding failures, and development diagnostics informed the final architecture and locally feasible backbone selection. Consequently, the controlled PLT results cannot be described as statistically independent held-out evidence of the complete design process.

The freeze boundary instead serves two purposes. First, after the final semantic/code revision, no prompt template, rubric, decision rule, model, target budget, search budget or source policy may be changed in response to final machine outcomes. Second, the preregistered external samples are selected and evaluated only after this freeze and are never used for model or prompt tuning.

This design sacrifices a stronger internal holdout claim in exchange for transparent reporting of the actual engineering history.

The definitive semantic/code freeze is commit
\texttt{286598d5eb642c3d632e7b756ccf1954fb922a13}; the final history-preserving repository snapshot on \texttt{main} is
\texttt{cc2de515cc481a6ec82c8c3edd5611209b84e97b}. Vericult uses \texttt{qwen3:4b} through Ollama at temperature 0. The exact installed model digest is captured by the runtime preflight before LIVE evidence acquisition.

The restored freeze was validated by GitHub Actions run \texttt{37608473802}. It incorporates the October 5 semantics previously validated at \texttt{c1944bb1df47a8668689c316e62563bdc700fbf3} but accidentally left on an unmerged branch.

The frozen configuration is:
\begin{itemize}
    \item prompt-level \texttt{cultural\_applicability\_v1} before D01--D10 planning;
    \item per-response \texttt{response\_assessability\_v1} before target extraction;
    \item at most three material targets per assessable response;
    \item two initial neutral verification questions per retrievable target;
    \item \texttt{top\_k}=3 accepted retrieval documents per query;
    \item at most two retrieval rounds, where round two contains at most one gap-specific follow-up per target;
    \item one semantic repair retry and one transport retry;
    \item deterministic span-ID normalization and Python-derived retrieval routing;
    \item per-dimension scorer recovery after persistent batch-structure failure;
    \item selective contextual fallback only for unresolved context-dependent recommendations;
    \item Tavily LIVE retrieval with advanced search depth;
    \item exclusion of the thesis repository and common annotation/result paths; and
    \item selective 0/1/2/\texttt{abstain} scoring with explicit \texttt{not\_culturally\_applicable}, \texttt{not\_assessable} and \texttt{insufficient\_evidence} outcomes.
\end{itemize}

The frozen verifier configuration has SHA-256
\texttt{3742f0ab4cb644760b7b455fdbd350ffbc6849c7579bc719749a57e70815c729}.

\begin{table}[htbp]
\centering
\footnotesize
\caption{Frozen controlled-experiment identifiers.}
\label{tab:frozen-experiment}
\begin{tabularx}{\textwidth}{p{4.2cm}X}
\toprule
\textbf{Component} & \textbf{Frozen value} \\
\midrule
Semantic/code revision & \texttt{286598d5eb642c3d632e7b756ccf1954fb922a13} \\
Repository freeze snapshot & \texttt{cc2de515cc481a6ec82c8c3edd5611209b84e97b} \\
Pipeline & \texttt{cultverify-1.1.0} \\
Backbone & Ollama \texttt{qwen3:4b}, temperature 0; exact runtime digest captured by preflight \\
PLT input & 30 prompts $\times$ 4 candidates; frozen SHA-256 above \\
Skywork baseline & \texttt{Skywork/Skywork-Reward-V2-Qwen3-4B}, revision \texttt{fd958fe} \\
Direct-judge order seed & \texttt{20260930} \\
Primary Vericult mode & REPLAY after audited LIVE evidence acquisition \\
\bottomrule
\end{tabularx}
\end{table}

The controlled experiment is deliberately split into evidence acquisition and evaluation. LIVE first runs all 30 Best-of-4 sets with the frozen configuration, allowing Tavily retrieval and writing evidence schedules, snapshots and candidate traces. The evidence audit requires 30 completed prompt records and 120 candidate traces, then hashes the LIVE output and all trace/evidence artifacts into a freeze manifest.

REPLAY verifies those hashes, reuses the same retrieval artifacts and performs no live search. REPLAY is the primary reported Vericult execution. Context planning, target comparison and scoring remain semantic model operations, so REPLAY freezes the evidence base rather than guaranteeing bitwise-identical model reasoning.

If a provider outage or execution-only failure occurs, the affected prompt may be rerun under the unchanged frozen code and configuration. Semantic changes are not permitted after freeze.

Vericult's primary task is single-response verification. The \texttt{rank} interface is an evaluation convenience requiring exactly four responses. Prompt context, cultural applicability and the D01--D10 plan are shared; response assessability and all downstream stages are evaluated per candidate.

The frozen ranking rule is selective and ordered:
\begin{enumerate}
    \item if any candidate fails technically, return \texttt{no\_clear\_winner};
    \item if the shared prompt is \texttt{not\_culturally\_applicable}, return that outcome rather than rank the responses;
    \item if any candidate is \texttt{not\_assessable}, return \texttt{not\_assessable};
    \item if any remaining candidate is \texttt{insufficient\_evidence}, return \texttt{insufficient\_evidence};
    \item exclude every \texttt{culturally\_inappropriate} candidate from endorsement;
    \item if all four candidates are excluded, return \texttt{no\_acceptable\_candidate};
    \item otherwise select the unique highest exact internal aggregate among eligible candidates; and
    \item if the eligible maximum is tied exactly, return \texttt{no\_clear\_winner}.
\end{enumerate}

There is no tie margin and no primary-dimension tie-break. Relative scores remain available diagnostically even when no candidate is endorsed. The result also records \texttt{coverage\_comparable}, indicating whether candidates were scored on the same set of non-abstained dimensions.

The frozen execution order is:
\begin{enumerate}
    \item run the local provider smoke gate and runtime preflight;
    \item acquire PLT evidence in LIVE mode;
    \item audit and hash the complete LIVE evidence/traces;
    \item run primary PLT REPLAY;
    \item run the Skywork reward-model baseline;
    \item run the direct LLM judge;
    \item only then join machine predictions to frozen Human Gold;
    \item evaluate the deterministic external-validation samples;
    \item compute the pre-declared metrics and paired statistics; and
    \item only after quantitative analysis select trace-backed cases for qualitative discussion.
\end{enumerate}

The final experiment runner rejects tracked semantic changes relative to the frozen code revision. Runtime artifacts are stored outside Git under \texttt{artifacts/final\_experiment/} and must be preserved with the thesis reproducibility archive.
